\documentclass[11pt]{article}
\usepackage[T1]{fontenc}
\usepackage{amsmath,amssymb,amsthm}
\usepackage[colorlinks=true,citecolor=blue,linkcolor=blue,urlcolor=blue]{hyperref}
\usepackage{enumitem}

\theoremstyle{plain}
\newtheorem{theorem}{Theorem}
\newtheorem{lemma}[theorem]{Lemma}
\newtheorem{corollary}[theorem]{Corollary}
\theoremstyle{definition}
\newtheorem{definition}[theorem]{Definition}
\theoremstyle{remark}
\newtheorem{remark}[theorem]{Remark}

\newcommand{\doi}[1]{\url{https://doi.org/#1}}
\newcommand{\Loc}{\mathrm{Loc}}
\newcommand{\NN}{\mathbb{N}}
\newcommand{\ZZ}{\mathbb{Z}}
\newcommand{\cM}{\mathcal{M}}

\title{A Prime-Generated Formalization of Nagata's Factoriality Theorem in Lean~4}
\author{Arthur Freitas Ramos}
\date{}

\begin{document}
\maketitle

\begin{abstract}
We present a Lean~4\,/\,Mathlib formalization of Nagata's factoriality theorem
for noetherian integral domains. The central hypothesis is a
\emph{prime-generated} condition on the multiplicative set: every element of the
submonoid is a finite product of elements that are both prime in the ring and
members of the submonoid. Under this hypothesis, if the localization is a
unique factorization domain (UFD), then so is the base ring. The formal
development comprises conditional transfer lemmas for divisibility,
irreducibility, and primality between a domain and a localization at a
prime-generated submonoid, together with the packaging of these ingredients into
a reusable Nagata theorem. As an application, we prove that if~$R$ is a
noetherian UFD then the polynomial ring~$R[X]$ is a UFD, by constructing a
Laurent-polynomial localization and invoking the formalized theorem. No public
formalization of Nagata's factoriality theorem is currently known to us from
available Lean, Coq, or Isabelle sources.
\end{abstract}

\section{Introduction}

Nagata's factoriality theorem is a standard tool in commutative algebra: it
recovers unique factorization in a base ring from unique factorization in a
suitable localization. The argument connects localization, irreducibility,
primality, and noetherianity in a way that makes it both a good test for a
formal algebra library and a useful building block for downstream results.

We describe a Lean~4 formalization, built on top of Mathlib, that packages a
complete proof of this theorem together with a worked application. An earlier
version of the project stated the main result under the hypothesis that every
element of the multiplicative set~$S$ is prime or a unit---a condition that is
too restrictive for a submonoid, since it essentially forces~$S$ to consist of
units and a single prime (up to associates). The present development replaces
that restricted statement with a \emph{prime-generated} hypothesis and rebuilds
the supporting localization machinery around the corrected formulation. The
earlier statement is retained only as a compatibility result.

\paragraph{Contributions.}
The formal contribution is threefold.
\begin{enumerate}[nosep]
  \item A formalization of Nagata's factoriality theorem under a prime-generated
  hypothesis on the multiplicative set~(\S\ref{sec:statement}).
  \item Conditional transfer lemmas that, under the prime-generated and
  prime-avoidance hypotheses arising in the Nagata argument, move divisibility,
  irreducibility, and primality information between a noetherian domain and its
  localization~(\S\ref{sec:architecture}).
  \item A theorem-driven application: if~$R$ is a noetherian UFD then~$R[X]$
  is a UFD, proved via a Laurent-polynomial localization and the formalized
  Nagata theorem~(\S\ref{sec:application}).
\end{enumerate}

\noindent
Section~\ref{sec:lessons} draws proof-engineering lessons from the
development; Section~\ref{sec:related} surveys related formal and mathematical
work; Section~\ref{sec:artifact} describes the artifact and its
reproducibility.


\section{Formal statement}\label{sec:statement}

We work over a commutative ring~$R$ that is both an integral domain and
noetherian.

\begin{definition}[Prime-generated submonoid]\label{def:primegen}
A submonoid~$S \subseteq R$ is \emph{prime-generated} if every
$s \in S$ can be written as a finite product
$s = q_1 \cdots q_n$
where each~$q_i$ satisfies $q_i \in S$ and $\mathrm{Prime}(q_i)$ in~$R$.
\end{definition}

\noindent
This is the right hypothesis for a localization-based factoriality argument: the
prime factors of denominators must themselves become units after localization, so
that denominators can be ``cleared'' factor by factor.

\begin{theorem}[Nagata]\label{thm:nagata}
Let $R$ be a noetherian integral domain and $S \subseteq R$ a prime-generated
submonoid. If the localization~$S^{-1}R$ is a UFD, then $R$~is a UFD.
\end{theorem}

The formal statement appears as \texttt{nagata\_theorem} in the file
\texttt{Nagata/Theorem.lean}. The earlier prime-or-unit variant is retained as
\texttt{nagata\_theorem\_of\_prime\_or\_unit}, but is proved separately as a
restricted compatibility theorem rather than derived from the prime-generated
statement.

\begin{remark}
The prime-or-unit condition, while convenient, is much more restrictive than it
appears. In a submonoid, if~$p$ and~$q$ are both non-unit primes in~$S$, then
$pq \in S$ but $pq$ is neither prime nor a unit. The prime-generated condition
avoids this collapse by asking only that every element of~$S$ \emph{factors}
into primes belonging to~$S$, rather than requiring each element to be prime
itself.
\end{remark}

\section{Proof architecture}\label{sec:architecture}

The proof proceeds in four layers, each providing reusable components for the
next.

\paragraph{1.\;Prime-generated submonoid API.}
The file \texttt{Nagata/Lemmas.lean} introduces the
predicate and establishes basic consequences: a prime-generated submonoid does
not contain zero; any prime element generates a prime-generated submonoid of its
powers; and products of prime-generated elements remain in the submonoid.

\paragraph{2.\;Conditional transfer lemmas.}
A series of lemmas transfers divisibility and factorization properties between
$R$ and $S^{-1}R$. These lemmas are \emph{conditional}: they rely on the
prime-generated hypothesis and on the prime-avoidance properties that it
entails. The key results are:
\begin{itemize}[nosep]
  \item An irreducible element of~$R$ that divides a product of primes from~$S$
    is itself prime.
  \item Divisibility in~$S^{-1}R$ can be lifted back to~$R$ when the
    denominators come from a prime-generated submonoid.
  \item Irreducibility and primality transfer between~$R$ and~$S^{-1}R$ under
    the prime-generated hypothesis.
\end{itemize}

\paragraph{3.\;Key lemma.}
The central technical result, formalized as
\texttt{nagata\_key\_lemma\_primeGenerated}, combines the transfer lemmas to
prove: if~$S$ is prime-generated and~$S^{-1}R$ is a UFD, then every irreducible
element of~$R$ is prime. This is the content that makes the Nagata argument
work.

\paragraph{4.\;Final packaging.}
The public theorem combines the key lemma with the standard fact that a
noetherian domain in which every irreducible is prime is a UFD.

\section{Application: polynomial rings}\label{sec:application}

The application file \texttt{Applications/Laurent.lean} puts the formalized
theorem to work.

\begin{corollary}\label{cor:poly}
If $R$ is a noetherian UFD, then the polynomial ring~$R[X]$ is a UFD.
\end{corollary}

The proof proceeds as follows. Let $S$ be the submonoid of $R[X]$ generated by
the powers of the indeterminate~$X$. Since $X$ is prime in $R[X]$, the
submonoid~$S$ is prime-generated. The localization $S^{-1}(R[X])$ is isomorphic
to the Laurent polynomial ring~$R[X, X^{-1}]$, which is a UFD when~$R$ is.
Theorem~\ref{thm:nagata} then gives unique factorization in~$R[X]$.

This application is not new mathematics, but it exercises the formalized Nagata
theorem in a nontrivial way: the proof genuinely depends on the prime-generated
localization machinery rather than reducing to a thin wrapper around a
pre-existing Mathlib instance.

\section{Lessons from the formalization}\label{sec:lessons}

Several aspects of the formal development may be of interest to the
proof-engineering community beyond the specific theorem.

\paragraph{Getting the hypothesis right.}
The most consequential design decision was the replacement of the
prime-or-unit hypothesis with the prime-generated one. The prime-or-unit
condition is superficially cleaner---it avoids quantifying over
multisets---but it is essentially degenerate for submonoids with more than one
non-unit prime generator. The formalization made this defect visible early,
because attempts to instantiate the theorem with natural examples (such as the
powers-of-$X$ submonoid) immediately failed.  The mathematical lesson is that a
condition on individual elements of a submonoid is the wrong abstraction level
when the argument requires factoring products of denominators; one must instead
ask for a \emph{factorization property} of the submonoid.

\paragraph{Transfer lemmas as conditional infrastructure.}
The divisibility, irreducibility, and primality transfer lemmas are not
unconditional equivalences between~$R$ and~$S^{-1}R$. They hold under the
specific combination of prime-generation and prime-avoidance hypotheses that
arises in the Nagata argument. Formalizing them as self-contained lemmas (rather
than inlining them into the main proof) required identifying exactly which
hypotheses each transfer result actually needs. This decomposition is
valuable not because each lemma is deep, but because the precise
statement of each lemma makes its reuse conditions explicit.

\paragraph{Automation and transparency.}
Lean's tactic automation handles much of the routine algebraic reasoning, but
the key proof obligations---particularly the case analysis on irreducible
elements dividing products of primes---require explicit combinatorial arguments
that cannot be hidden behind \texttt{simp} or \texttt{ring}. This is a feature,
not a limitation: the formal proof preserves exactly the case distinctions that
constitute the mathematical content of the argument.

\paragraph{Localization ergonomics.}
Working with Mathlib's localization API involves navigating between the abstract
\texttt{IsLocalization} typeclass and concrete \texttt{Localization} instances.
The formal development stays close to the \texttt{Localization} type for
definitional convenience, which simplifies denominator manipulation at the cost
of some generality. A future Mathlib-oriented refactoring might abstract over
the localization construction, but for the present theorem the concrete approach
proved more practical.

\section{Related work}\label{sec:related}

\paragraph{Mathematical background.}
The classical result appears in Nagata's treatment of localization and
factoriality~\cite{nagata1962}; further context is provided by
Samuel~\cite{samuel1964} and Matsumura~\cite{matsumura1989}.

\paragraph{Formal algebra libraries.}
The Lean baseline for this project is Mathlib~\cite{mathlib}, which provides the
commutative algebra, localization, and UFD infrastructure on which our
development builds. Our contribution is not the formalization of localization or
UFD machinery, but the formal packaging of a specific Nagata-style theorem and
the proof architecture needed to make it applicable.

In Coq, the Mathematical Components library~\cite{mathcomp} offers a mature
algebraic environment with strong support for divisibility and polynomial
algebra. In Isabelle/HOL, Bordg~\cite{bordg-localization} formalizes
localization of commutative rings in the Archive of Formal Proofs, and
Divas{\'o}n et al.~\cite{divason-berlekamp} formalize the
Berlekamp--Zassenhaus factorization algorithm, demonstrating substantial
factorization infrastructure in that system.

\paragraph{Novelty.}
No public formalization of Nagata's factoriality theorem is currently known to
us from available Lean/Mathlib, Coq/MathComp, or Isabelle/AFP sources. This is
a statement about the present public state of these libraries, not a claim that
no unpublished or in-progress formalization exists.

\section{Prospective library contributions}

Not every result in the development is a candidate for upstreaming.
\begin{itemize}[nosep]
  \item \textbf{Good candidates.} The localization helpers and
  generic submonoid lemmas (e.g., that a prime-generated submonoid does not
  contain zero) are reusable beyond the present theorem.
  \item \textbf{Conditional candidates.} The prime-generated interface and its
  factorization lemmas are mathematically natural, but would need API
  stabilization before they fit Mathlib's abstraction boundaries.
  \item \textbf{Paper-level contribution.} The Nagata theorem itself and its
  specific proof packaging are the core contribution of this repository.
\end{itemize}

\section{Artifact and reproducibility}\label{sec:artifact}

The development targets Lean~4.24.0, builds with Lake, and contains no
\texttt{sorry}, \texttt{admit}, or \texttt{axiom} placeholders.
The repository comprises 16~Lean source files totaling approximately
925~lines of Lean under \texttt{NagataFactoriality/}, with 84~theorem
and lemma declarations. Continuous integration and citation metadata
(\texttt{CITATION.cff}) are included.

\paragraph{Building the artifact.}
\begin{verbatim}
lake exe cache get
lake build
\end{verbatim}

\paragraph{Release packaging.}
The submission artifact is a tagged source release that freezes the
\texttt{lean-toolchain}, \texttt{lake-manifest.json}, the Lean sources, and
this paper. The paper cites that tagged release (or, at minimum, a specific
commit hash) rather than a moving branch. The same release can be mirrored to a
DOI-minting archive such as Zenodo if the venue accepts archival artifacts.

\section{Limitations and future work}

The formalization covers the multiplicative version of Nagata's theorem (for
submonoids generated by primes). The more general ideal-theoretic formulations
of the result are not addressed. On the application side, the Laurent-polynomial
example is intentionally compact; richer applications (e.g., to power-series
rings or to Dedekind domains) would further exercise the formalized machinery.
Finally, some localization API choices are tailored to this project and may need
refactoring for eventual Mathlib integration.

\begin{thebibliography}{9}

\bibitem{bordg-localization}
A.~Bordg.
\newblock The localization of a commutative ring.
\newblock \emph{Archive of Formal Proofs}, June 2018.
\newblock \url{https://www.isa-afp.org/entries/Localization_Ring.html}.

\bibitem{divason-berlekamp}
J.~Divas{\'o}n, S.\,J.\,C.~Joosten, R.~Thiemann, and A.~Yamada.
\newblock A verified implementation of the {B}erlekamp--{Z}assenhaus
  factorization algorithm.
\newblock \emph{J.~Automated Reasoning}, 64:699--735, 2020.
\newblock \doi{10.1007/s10817-019-09526-y}.

\bibitem{mathlib}
The mathlib Community.
\newblock The {L}ean mathematical library.
\newblock In \emph{Proc.\ 9th ACM SIGPLAN Int.\ Conf.\ Certified Programs and
  Proofs (CPP~2020)}, pages 367--381, 2020.
\newblock \doi{10.1145/3372885.3373824}.

\bibitem{mathcomp}
A.~Mahboubi and E.~Tassi.
\newblock \emph{Mathematical Components}.
\newblock Zenodo, 2022.
\newblock \doi{10.5281/zenodo.3999478}.

\bibitem{matsumura1989}
H.~Matsumura.
\newblock \emph{Commutative Ring Theory}.
\newblock Cambridge University Press, 1989.

\bibitem{nagata1962}
M.~Nagata.
\newblock \emph{Local Rings}.
\newblock Interscience, 1962.

\bibitem{samuel1964}
P.~Samuel.
\newblock \emph{Lectures on Unique Factorization Domains}.
\newblock Tata Institute of Fundamental Research, 1964.

\end{thebibliography}

\end{document}
